\section{学习包的结构与用法}

\subsection{文件夹}
学习包的文件夹名为 \Verb|应用线性回归_R实操/|。在本机，它放在 \Verb|应用线性回归/| 下，与原有的 \Verb|回归1|–\Verb|回归4|、\Verb|hw1|–\Verb|hw4| 并列，原文件夹没有改动；从学习网站下载的压缩包解压后也是这个文件夹。结构如下：

\begin{ROut}[目录结构]
应用线性回归_R实操/
|-- 应用线性回归_R实操.Rproj   ← 双击它打开 RStudio
|-- 00_安装与检查.R            ← 第一次打开时运行一次
|-- README.md
|-- data/                      ← 全部数据（24 个 CSV）
|   |-- 数据说明.md            ← 每个文件的来源、变量、在哪里用到
|   `-- make_inline_data.R     ← 把课件里写在代码中的数据导出为 CSV
|-- scripts/                   ← 按讲次整理的可运行脚本（本笔记第 1–8 章、附录 B）
|   |-- 01_R工作流与数据处理.R
|   |-- 02_多重回归_拟合与解读.R
|   |-- 03_假设检验与变量筛选.R
|   |-- 04_回归诊断.R
|   |-- 05_共线性与岭回归.R
|   |-- 06_一般线性模型_分类变量与方差分析.R
|   |-- 07_协方差分析.R
|   |-- 08_重复测量方差分析.R
|   `-- 09_课件笔记例题复现.R
|-- homework/                  ← 四次作业的代码（本笔记第 9 章）
|-- 课堂原始代码/               ← 老师发的原件，未改动；附对照说明.md
|-- output/                    ← 运行脚本时生成的图和表
`-- 笔记/                      ← 本笔记的 PDF 与 LaTeX 源文件
\end{ROut}

\subsection{打开与运行}
\begin{enumerate}
\item 双击 \Verb|应用线性回归_R实操.Rproj|。RStudio 会把工作目录设为学习包根目录，所有脚本里的 \Verb|read.csv("data/...")| 都能直接找到文件。不要用“直接双击 .R 文件”的方式打开，那样工作目录不对。
\item 第一次使用先运行 \Verb|00_安装与检查.R|，它会装好缺少的包（car、leaps、olsrr、ppcor、lmtest、sandwich、emmeans、afex、ez、tidyr、ggplot2），并检查数据是否齐全。
\item 打开 \Verb|scripts/| 下的脚本，按 Ctrl/Cmd + Shift + O 显示右侧大纲。脚本里形如 \Verb|# 2.7 预测：均数的 CI 与个体的 PI ----| 的行就是大纲条目。脚本编号与本笔记的章号相同（脚本 02 对应第 2 章），正文提到“脚本 §2.7”时，指的就是大纲里的这一条。
\item 光标放在某一行按 Ctrl/Cmd + Enter 运行这一行（或选中的几行）；Ctrl/Cmd + Shift + Enter 从头运行整份脚本。每份脚本都已从头到尾跑通。
\end{enumerate}

\subsection{怎样配合本笔记使用}
每章开头的小字注明对应的课件页码、脚本和数据。正文只摘出关键的几行代码和输出，完整代码在脚本里。建议的顺序是：先读本章正文弄清“要做什么、读哪几个数”，再在 RStudio 里逐段运行脚本，对照正文核对数值，最后打开对应的作业脚本看同样的方法怎么用在作业上。

\begin{table}[htbp]
\centering\small
\caption{课件、脚本、数据与作业的对应关系}
\label{tab:map}
\begin{tabular}{@{}p{3.3cm}p{4.2cm}p{4.3cm}p{2.0cm}@{}}
\toprule
课件 & 脚本 & 主要数据 & 对应作业\\
\midrule
—— & 01 R 工作流与数据处理 & 各数据 & ——\\
第一讲 模型与解释 & 02 多重回归：拟合与解读 & datasetA & 作业一\\
第一讲后半 变量筛选 & 03 假设检验与变量筛选 & datasetC、diabetes27 & 作业一\\
第二讲 回归诊断 & 04 回归诊断 & datasetB & 作业二第一题\\
第二讲 共线性与岭回归 & 05 共线性与岭回归 & datasetB & 作业二第二题\\
第三讲 一般线性模型 & 06 一般线性模型 & birthwt、glucose\_block、sbp\_factorial2x2 & 作业三\\
第三、四讲 协方差分析 & 07 协方差分析 & datasetC\_glm、bp、knee & 作业四第 1 题\\
第四讲 重复测量 & 08 重复测量方差分析 & dm、df\_long & 作业四第 2 题\\
课件笔记各章例题 & 09 课件笔记例题复现 & 12 个小数据集 & ——\\
\bottomrule
\end{tabular}
\TableNote{数据文件均在 \texttt{data/} 下；文件来源与变量见 \texttt{data/数据说明.md}。}
\end{table}

\subsection{网页版与下载}
这份笔记也收进了学习网站《医学统计学 · 学习笔记》，是其中的第四门课程“应用线性回归 · R 实操”：

\begin{center}
\url{https://fxt-gw-pb.github.io/public-health-statistics-notes/#/course/rcode}
\end{center}

网页版可以在线阅读、全文搜索、一键复制代码块；课程页提供两个下载入口：本笔记的 PDF，以及完整的 R 学习包（\Verb|.Rproj|、全部数据、九份讲次脚本、四份作业脚本、课堂原始代码和本笔记的 LaTeX 源文件）。下载学习包解压后，按上一节的步骤打开即可。

\begin{notebox}[关于课堂原始代码]
老师发的脚本保存在 \Verb|课堂原始代码/| 里，没有改动。它们写的是 \Verb|read.csv("d://data/...")| 这样的绝对路径，又大量使用 \Verb|attach()|，有几处直接运行会报错（例如第四讲读入的对象叫 \Verb|da| 后面却用 \Verb|bp|）。整理版把这些问题都改掉了，逐条对照见 \Verb|课堂原始代码/对照说明.md| 和附录 D。
\end{notebox}
